% 切丛与余切丛 (Tangent and Cotangent Bundles)
% 描述: 展示切丛和余切丛的结构，以及它们之间的对偶关系
% 数学概念: 切丛、余切丛、向量场、微分形式、李括号
\documentclass[border=10pt]{standalone}
\usepackage{tikz}
\usepackage{amsmath,amssymb}
\usetikzlibrary{arrows.meta,positioning,calc,patterns,decorations.pathmorphing}

\begin{document}
\begin{tikzpicture}[
    >=Stealth,
    scale=1.0
]

% ============= 左侧: 切丛 =============
\begin{scope}[xshift=-5cm, yshift=1cm]
    % 标题
    \node[font=\bfseries] at (0, 3) {切丛 $TM$};
    
    % 底流形M
    \draw[color=blue!60!black, line width=1pt, fill=blue!10, opacity=0.5] 
        (-1.5, -1.8) .. controls (-0.5, -2.2) and (0.5, -1.8) .. (1.5, -1.8);
    \node[color=blue!60!black, font=\small] at (1.7, -1.8) {$M$};
    
    % 点p
    \fill[black] (0, -1.8) circle (2pt);
    \node[below] at (0, -1.8) {$p$};
    
    % 切空间纤维
    \draw[color=red!70!black, line width=1pt, fill=red!10, opacity=0.3] 
        (-0.8, -0.5) -- (-0.8, 2) -- (0.8, 2) -- (0.8, -0.5) -- cycle;
    \draw[color=red!70!black, line width=1.5pt, dashed] (-0.8, -0.5) -- (-0.8, 2);
    \draw[color=red!70!black, line width=1.5pt, dashed] (0.8, -0.5) -- (0.8, 2);
    \node[color=red!70!black, font=\small, right] at (0.9, 0.8) {$T_pM$};
    
    % 切向量
    \draw[->, color=green!60!black, line width=1.5pt] (0, 0.5) -- (0.5, 1.5);
    \fill[black] (0, 0.5) circle (2pt);
    \node[color=green!60!black, font=\small, right] at (0.3, 1.1) {$v$};
    \node[right, font=\tiny] at (0.1, 0.3) {$(p, v)$};
    
    % 投影
    \draw[->, color=orange!80!black, line width=1pt] (0, -0.2) -- (0, -1.5);
    \node[color=orange!80!black, font=\tiny, right] at (0.1, -0.8) {$\pi$};
\end{scope}

% ============= 中间: 余切丛 =============
\begin{scope}[xshift=0cm, yshift=1cm]
    % 标题
    \node[font=\bfseries] at (0, 3) {余切丛 $T^*M$};
    
    % 底流形M
    \draw[color=blue!60!black, line width=1pt, fill=blue!10, opacity=0.5] 
        (-1.5, -1.8) .. controls (-0.5, -2.2) and (0.5, -1.8) .. (1.5, -1.8);
    \node[color=blue!60!black, font=\small] at (1.7, -1.8) {$M$};
    
    % 点p
    \fill[black] (0, -1.8) circle (2pt);
    \node[below] at (0, -1.8) {$p$};
    
    % 余切空间纤维
    \draw[color=purple!70!black, line width=1pt, fill=purple!10, opacity=0.3] 
        (-0.8, -0.5) -- (-0.8, 2) -- (0.8, 2) -- (0.8, -0.5) -- cycle;
    \draw[color=purple!70!black, line width=1.5pt, dashed] (-0.8, -0.5) -- (-0.8, 2);
    \draw[color=purple!70!black, line width=1.5pt, dashed] (0.8, -0.5) -- (0.8, 2);
    \node[color=purple!70!black, font=\small, right] at (0.9, 0.8) {$T^*_pM$};
    
    % 余切向量(1-形式)
    \draw[color=cyan!70!black, line width=2pt] (-0.3, 0.8) -- (0.3, 0.8);
    \draw[color=cyan!70!black, line width=1.5pt, ->] (-0.3, 0.8) -- (-0.3, 1.3);
    \draw[color=cyan!70!black, line width=1.5pt, ->] (0.3, 0.8) -- (0.3, 1.3);
    \fill[black] (0, 0.5) circle (2pt);
    \node[color=cyan!70!black, font=\small, right] at (0.4, 1.1) {$\alpha$};
    
    % 投影
    \draw[->, color=orange!80!black, line width=1pt] (0, -0.2) -- (0, -1.5);
    \node[color=orange!80!black, font=\tiny, right] at (0.1, -0.8) {$\pi^*$};
\end{scope}

% ============= 右侧: 对偶配对 =============
\begin{scope}[xshift=5cm, yshift=1cm]
    % 标题
    \node[font=\bfseries] at (0, 3) {对偶配对};
    
    % 两个对偶空间
    % TM
    \draw[color=red!70!black, line width=1pt, fill=red!10, opacity=0.3] 
        (-1.2, -0.5) -- (-1.2, 1.8) -- (-0.2, 1.8) -- (-0.2, -0.5) -- cycle;
    \node[color=red!70!black, font=\small, left] at (-1.3, 0.8) {$T_pM$};
    
    % T*M
    \draw[color=purple!70!black, line width=1pt, fill=purple!10, opacity=0.3] 
        (0.2, -0.5) -- (0.2, 1.8) -- (1.2, 1.8) -- (1.2, -0.5) -- cycle;
    \node[color=purple!70!black, font=\small, right] at (1.3, 0.8) {$T^*_pM$};
    
    % 向量
    \draw[->, color=green!60!black, line width=1.5pt] (-0.7, 0.5) -- (-0.2, 1.2);
    \node[color=green!60!black, font=\small] at (-0.6, 1) {$v$};
    
    % 1-形式
    \draw[color=cyan!70!black, line width=1.5pt] (0.5, 0.8) -- (0.9, 0.8);
    \draw[color=cyan!70!black, line width=1pt, ->] (0.5, 0.8) -- (0.5, 1.2);
    \draw[color=cyan!70!black, line width=1pt, ->] (0.9, 0.8) -- (0.9, 1.2);
    \node[color=cyan!70!black, font=\small] at (0.9, 1.1) {$\alpha$};
    
    % 配对值
    \node[font=\Large] at (0, -0.8) {$\langle \alpha, v \rangle \in \mathbb{R}$};
    
    % 说明
    \node[anchor=north, font=\scriptsize, text width=4cm, align=center] 
        at (0, -1.5) {线性泛函作用\\
                       $\alpha(v) = \alpha_i v^i$};
\end{scope}

% ============= 底部: 向量场与李括号 =============
\node[anchor=north, font=\small, draw=blue!30, fill=blue!5, rounded corners, inner sep=10pt] 
    at (4, -4.5) {
    \begin{tabular}{ll}
        \textbf{向量场空间:} & \\
        $\mathfrak{X}(M) = \Gamma(TM)$: & 光滑向量场 \\
        $\Omega^1(M) = \Gamma(T^*M)$: & 光滑1-形式 \\[5pt]
        \textbf{李括号:} & $[X, Y] \in \mathfrak{X}(M)$ \\
        定义: & $[X, Y](f) = X(Y(f)) - Y(X(f))$ \\[5pt]
        \textbf{结构:} & $(\mathfrak{X}(M), [\cdot, \cdot])$ 构成Lie代数
    \end{tabular}
};

% ============= 局部坐标 =============
\begin{scope}[xshift=-5cm, yshift=-4cm]
    \node[anchor=north west, font=\scriptsize, text width=5cm, align=left] 
        at (0, 0) {
        \textbf{局部坐标:}\\[5pt]
        切向量: $v = v^i \frac{\partial}{\partial x^i}$\\[5pt]
        1-形式: $\alpha = \alpha_i dx^i$\\[5pt]
        对偶基: $\langle dx^i, \frac{\partial}{\partial x^j} \rangle = \delta^i_j$
    };
\end{scope}

% ============= 辛结构 =============
\begin{scope}[xshift=5cm, yshift=-4cm]
    \node[anchor=north east, font=\scriptsize, text width=5cm, align=left] 
        at (0, 0) {
        \textbf{余切丛的辛结构:}\\[5pt]
        典范1-形式:\\
        $\theta_{(p,\alpha)}(V) = \alpha(\pi_* V)$\\[5pt]
        辛形式:\\
        $\omega = -d\theta$\\
        $(T^*M, \omega)$是辛流形
    };
\end{scope}

\end{tikzpicture}
\end{document}
